%% ============================================================================
%%  Two alternative framings of the Results, to compare. Not \input anywhere;
%%  copy the preferred one into main.tex. Numbers are from the current fits.
%% ============================================================================


%% ============================================================================
%%  VERSION A - ecological framing: supply-limited vs habitat-limited
%%  Story: larvae reach everywhere; habitat filters where cockles occur; local
%%  (spatial) processes set how many; larval supply adds nothing and is, if
%%  anything, highest where habitat is poor.
%% ============================================================================

\section{Results}

\subsection{Cockles occur in a few dense beds, although larvae reach the whole fjord}

Cockles were present at 1{,}527 of 8{,}137 stations (18.8\%), and biomass was concentrated in a few dense beds along the central channel and in the western and southern basins (Fig.~\ref{fig:predictions}). This patchiness was not matched by larval supply: simulated larvae dispersed widely from every release cell (Fig.~\ref{fig:study_area}b), and every survey station received some larval input (presence-weighted in-strength $>0$ at all stations). Larvae therefore reached areas both with and without established beds.

\subsection{Habitat determines where cockles occur}

Occurrence was strongly structured by the environment (Fig.~\ref{fig:coefficients}). The probability of presence peaked at intermediate depths, with a maximum at 7.1~m (95\% CI 5.9--8.0~m; depth$^2$: $\hat{\beta}=-0.30$, $-0.43$ to $-0.17$ per SD), increased with salinity ($0.82$, $0.43$ to $1.21$), and declined with temperature ($-0.51$, $-0.89$ to $-0.13$) and maximum shear stress ($-0.68$, $-0.99$ to $-0.36$). Adding the environment to a model containing only the spatial field improved both the fit ($\Delta$AIC $=-59$) and the prediction of held-out grid cells ($\Delta$ELPD $=86.2$, SE $16.4$; Table~\ref{tab:model-selection}).

\subsection{Local processes determine how many}

Where cockles were present, their biomass was largely unrelated to the environment: only shear stress ($-0.31$, $-0.61$ to $-0.01$) and the number of days of low oxygen ($-0.59$, $-1.14$ to $-0.05$) had CIs excluding zero (Fig.~\ref{fig:coefficients}). Instead, biomass and occurrence both showed strong residual spatial structure, decaying over a Mat\'ern range of 2.7~km (95\% CI 1.8--4.0~km) and 2.2~km (1.8--2.8~km), respectively, i.e.\ over about one 2~$\times$~2~km grid cell. Beds were thus patchier than the environmental gradients, at scales at or below the resolution of the connectivity model.

\subsection{Larval supply does not explain cockle distribution}

Larval connectivity added no explanatory or predictive power. Adding in-strength to the Space + Environment model did not improve the fit ($\Delta$AIC $=0.1$) and slightly reduced out-of-sample predictive performance ($\Delta$ELPD $=-20.7$, SE $12.6$); added to the spatial field alone, it reduced it further ($-36.0$, SE $10.9$; Table~\ref{tab:model-selection}). Its coefficient was negative in both components, with CIs including zero (presence: $-0.27$, $-0.55$ to $0.00$; biomass where present: $-0.04$, $-0.27$ to $0.20$), and predicted presence and biomass changed little across its observed range (Fig.~\ref{fig:coefficients}b). The negative tendency reflects where larvae accumulate: in-strength was highest in deeper ($\rho=0.50$), less saline ($\rho=-0.55$) and more often hypoxic ($\rho=0.37$) areas (Fig.~\ref{fig:connectivity_environment}), i.e.\ where habitat for cockles was poor. The same conclusions held for the four alternative connectivity metrics, none of which improved the fit ($\Delta$AIC $+1.7$ to $+3.6$ relative to in-strength; Table~\ref{tab:connectivity-metrics}) or had a coefficient with a CI excluding zero (Fig.~\ref{fig:connectivity_coefficients}).

\begin{table}[ht]
\centering
\small
\caption{Model selection for cockle presence, biomass where present and both together (the delta-gamma model). $\Delta$AIC is the difference in AIC to the best model (lower AIC is better). $\Delta$ELPD is the difference in expected log predictive density to the best model in spatially blocked 10-fold cross-validation (higher ELPD is better), with its standard error in parentheses. The best model in each column is shown in bold. The total AIC and ELPD are the sums of the two parts. All models include survey (gear) and year as factors and a spatial random field. Connectivity is presence-weighted in-strength; environment comprises depth (linear and quadratic), temperature, oxygen, salinity, and maximum shear stress. A dash marks a model that did not converge.}
\label{tab:model-selection}
\begin{tabular}{lrrrrrr}
\toprule
 & \multicolumn{3}{c}{$\Delta$AIC} & \multicolumn{3}{c}{$\Delta$ELPD (SE)} \\
\cmidrule(lr){2-4} \cmidrule(lr){5-7}
Model & Presence & Biomass & Total & Presence & Biomass & Total \\
\midrule
Space + Environment & 1.8 & 1.6 & \textbf{0.0} & -18.8 (4.6) & -51.9 (32.9) & -33.7 (24.7) \\
Space + Environment + Connectivity & \textbf{0.0} & 3.5 & 0.1 & \textbf{0.0} & -37.0 (10.9) & \textbf{0.0} \\
Space & 62.8 & \textbf{0.0} & 59.3 & -87.2 (10.9) & \textbf{0.0} & -50.2 (15.6) \\
Space + Connectivity & 63.8 & 1.9 & 62.2 & -93.9 (10.7) & -50.4 (29.7) & -107.2 (24.9) \\
\bottomrule
\end{tabular}
\end{table}




%% ============================================================================
%%  VERSION B - methodological framing: testing connectivity in SDMs
%%  Story: a structured test - does connectivity add to an SDM with
%%  environment and space? First the model comparison (the test), then the
%%  effects that explain the outcome, then robustness.
%% ============================================================================

\section{Results}

\subsection{Connectivity did not improve the species distribution models}

Of the four model structures, Space + Environment was best supported, both in sample and out of sample (Table~\ref{tab:model-selection}). Adding presence-weighted in-strength to it did not improve the fit ($\Delta$AIC $=0.1$) and did not improve the prediction of held-out grid cells ($\Delta$ELPD $=-20.7$, SE $12.6$). In contrast, the environmental predictors improved both: removing them from the full model increased AIC by $\approx$62, and adding them to the spatial field increased ELPD by $86.2$ (SE $16.4$), and by $101.1$ (SE $17.7$) when in-strength was also included. A model with in-strength but no environment performed worst ($\Delta$AIC $=62.2$; $\Delta$ELPD $=-122.2$, SE $15.9$, relative to the best model).

The result did not depend on how connectivity was summarised. Replacing in-strength by in-degree, in-closeness, eigenvector centrality or transitivity did not improve the fit of the full model ($\Delta$AIC $+1.7$ to $+3.6$; Table~\ref{tab:connectivity-metrics}), and none of the five metrics had a coefficient whose CI excluded zero in either component (Fig.~\ref{fig:connectivity_coefficients}).

\begin{table}[ht]
\centering
\small
\caption{Model selection for cockle presence, biomass where present and both together (the delta-gamma model). $\Delta$AIC is the difference in AIC to the best model (lower AIC is better). $\Delta$ELPD is the difference in expected log predictive density to the best model in spatially blocked 10-fold cross-validation (higher ELPD is better), with its standard error in parentheses. The best model in each column is shown in bold. The total AIC and ELPD are the sums of the two parts. All models include survey (gear) and year as factors and a spatial random field. Connectivity is presence-weighted in-strength; environment comprises depth (linear and quadratic), temperature, oxygen, salinity, and maximum shear stress. A dash marks a model that did not converge.}
\label{tab:model-selection}
\begin{tabular}{lrrrrrr}
\toprule
 & \multicolumn{3}{c}{$\Delta$AIC} & \multicolumn{3}{c}{$\Delta$ELPD (SE)} \\
\cmidrule(lr){2-4} \cmidrule(lr){5-7}
Model & Presence & Biomass & Total & Presence & Biomass & Total \\
\midrule
Space + Environment & 1.8 & 1.6 & \textbf{0.0} & -18.8 (4.6) & -51.9 (32.9) & -33.7 (24.7) \\
Space + Environment + Connectivity & \textbf{0.0} & 3.5 & 0.1 & \textbf{0.0} & -37.0 (10.9) & \textbf{0.0} \\
Space & 62.8 & \textbf{0.0} & 59.3 & -87.2 (10.9) & \textbf{0.0} & -50.2 (15.6) \\
Space + Connectivity & 63.8 & 1.9 & 62.2 & -93.9 (10.7) & -50.4 (29.7) & -107.2 (24.9) \\
\bottomrule
\end{tabular}
\end{table}


\subsection{Environmental and connectivity effects}

The environment acted mainly on occurrence (Fig.~\ref{fig:coefficients}). The probability of presence peaked at 7.1~m depth (95\% CI 5.9--8.0~m; depth$^2$: $\hat{\beta}=-0.30$, $-0.43$ to $-0.17$ per SD), increased with salinity ($0.82$, $0.43$ to $1.21$), and declined with temperature ($-0.51$, $-0.89$ to $-0.13$) and maximum shear stress ($-0.68$, $-0.99$ to $-0.36$). Biomass where present declined with shear stress ($-0.31$, $-0.61$ to $-0.01$) and the number of days of low oxygen ($-0.59$, $-1.14$ to $-0.05$), whereas the other environmental effects on biomass had CIs including zero.

The in-strength coefficient was negative in both components, with CIs including zero (presence: $-0.27$, $-0.55$ to $0.00$; biomass where present: $-0.04$, $-0.27$ to $0.20$), and predicted presence and biomass changed little across its observed range (Fig.~\ref{fig:coefficients}b). In-strength was correlated with the environment: it was higher in deeper ($\rho=0.50$), less saline ($\rho=-0.55$), less exposed ($\rho=-0.52$ with shear stress) and more often hypoxic ($\rho=0.37$) areas (Fig.~\ref{fig:connectivity_environment}).

\subsection{Spatial structure}

Residual spatial autocorrelation decayed over a Mat\'ern range of 2.2~km (95\% CI 1.8--2.8~km) for presence and 2.7~km (1.8--4.0~km) for biomass where present, i.e.\ over about one grid cell of the connectivity model. Predicted biomass was concentrated in a few dense beds along the central channel and in the western and southern basins, with prediction uncertainty lowest in the densely surveyed beds and highest in the sparsely surveyed north-eastern basin and northern arm (Fig.~\ref{fig:predictions}).
